\scanpage{248}

% top, heading and small number above it
\mspage{1}Relations des dérivateurs avec la catégorie Hot.

% lines 2–3
Soit $\mathbb D$ un dérivateur, $\mathcal A$ sa catégorie fondamentale

% line 3, centred
\[\mathcal A=\mathbb D(e)\]

% lines 4–5
Je me propose de définir, \deleted{\unsure{dès}} \unsure{si} faire se peut, des foncteurs

% upper display (1), under one left brace, with right-hand formulas
\[
\left\{\begin{aligned}
\mathrm{Hot}\times\mathcal A&\xrightarrow{\pi_{\mathbb D}}\mathcal A
&(\unsure{\xi|\zeta}\,\unsure{,|\times}\,\deleted{\unsure{F|M}}\unsure{M|F})&\longrightarrow\xi\times F\ \text{ou }\xi\otimes\deleted{\unsure{F}}M\\
&&&\text{(suivant les cas)}\\
\mathrm{Hot}^{\circ}\times\mathcal A&\xrightarrow{\varpi_{\mathbb D}}\mathcal A
&(\xi,\deleted{\unsure{F}}\unsure{M|A})&\longrightarrow
\unsure{\mathrm{hom}|\mathrm{Hom}}(\xi,\deleted{\unsure{F}}M)
\end{aligned}\right.\tag*{(1)}
\]

% below (1), lines 1–3
Ces foncteurs doivent s'échanger par dualité : si on prend les foncteurs $\pi^\circ$, $\varpi^\circ$ sur les catégories opposées, on trouve

% middle display (1′)
\[
\left\{\begin{aligned}
\unsure{\mathrm{HOT}|\mathrm{Hot}}\times\mathcal A^\circ&\xrightarrow{\deleted{\varpi^\circ}\ \varpi_{\mathbb D}^{\circ}}\mathcal A^\circ\\
\mathrm{Hot}^{\circ}\times\mathcal A&\xrightarrow{\pi_{\mathbb D}^{\circ}}\mathcal A^\circ
\end{aligned}\right.\tag*{(\uline{1}')}
\]

% above duality identities
et on doit avoir

% middle, two identities and struck label
\[
\left\{\begin{aligned}
\varpi_{\mathbb D}^\circ&=\pi_{\mathbb D^\circ}\\
\deleted{\unsure{\varpi|\pi}}\pi_{\mathbb D}^\circ&=\varpi_{\mathbb D^\circ}
\end{aligned}\right.\ .\tag*{\deleted{\unsure{(2)}}}
\]

% line immediately below the identities
\deleted{\unsure{Les|Leurs} propriétés fondamentales}

% lower half, introduction to functor form
À vrai dire, considérant (1) comme des foncteurs

% lower half, first unnumbered pair
\[\begin{aligned}
\mathrm{Hot}&\xrightarrow{\pi_{\mathbb D}}\underline{\mathrm{Hom}}(\mathcal A,\mathcal A)\\
\mathrm{Hot}^{\circ}&\xrightarrow{\varpi_{\mathbb D}}\underline{\mathrm{Hom}}(\mathcal A,\mathcal A)
\end{aligned}\]

% before (2)
on s'attend à ce que ces foncteurs s'étendent en

% lower display (2)
\[
\left\{\begin{aligned}
\mathrm{Hot}&\xrightarrow{\pi_{\mathbb D}}\underline{\mathrm{Hom}}(\mathbb D,\mathbb D)\\
\mathrm{Hot}^{\circ}&\xrightarrow{\varpi_{\mathbb D}}\underline{\mathrm{Hom}}(\mathcal A,\mathcal A)
\end{aligned}\right.\tag*{(2)}
\]

% foot, last four lines with added qualification
Cela signifie, essentiellement, que \inserted{pour $\xi\in\mathrm{Ob}\,\mathrm{Hot}$ donné,} les opérations $\xi\times_{\mathbb D}M$ et $\underline{\mathrm{Hom}}_{\mathbb D}(\xi,F)$ sont définies non seulement dans $\mathcal A=\mathbb D(e)$, mais dans tous les $\mathbb D(I)$, et ceci de façon à commuter aux foncteurs $f^*$ :

% archive sheet number: 248
